%%%%%%%%%%%%%%%%%%%%%%%%%%%%%%%%%%%%%%%%%
% LAGG curriculum vitae
%
% The original template (Wilson Resume/CV) has been downloaded from:
% http://www.LaTeXTemplates.com
%
% Original authors:
% Howard Wilson (https://github.com/watsonbox/cv_template_2004) with
% extensive modifications by Vel (vel@latextemplates.com)
%
% License:
% CC BY-NC-SA 3.0 (http://creativecommons.org/licenses/by-nc-sa/3.0/)
%
%%%%%%%%%%%%%%%%%%%%%%%%%%%%%%%%%%%%%%%%%

%----------------------------------------------------------------------------------------
%	PACKAGES AND OTHER DOCUMENT CONFIGURATIONS
%----------------------------------------------------------------------------------------

\documentclass[10pt]{article} % Default font size

%%%%%%%%%%%%%%%%%%%%%%%%%%%%%%%%%%%%%%%%%
% Structure specification file
%
% The original file (Wilson Resume/CV) has been downloaded from:
% http://www.LaTeXTemplates.com
%
% License:
% CC BY-NC-SA 3.0 (http://creativecommons.org/licenses/by-nc-sa/3.0/)
%
%%%%%%%%%%%%%%%%%%%%%%%%%%%%%%%%%%%%%%%%%

%----------------------------------------------------------------------------------------
%	PACKAGES AND OTHER DOCUMENT CONFIGURATIONS
%----------------------------------------------------------------------------------------

% Use A4 paper and set margins
\usepackage[a4paper, hmargin=25mm, vmargin=30mm, top=20mm]{geometry}

% Use the Erewhon font
%\usepackage[proportional,scaled=1.064]{erewhon}
%\usepackage[erewhon,vvarbb,bigdelims]{newtxmath}

% Required for inputting international characters
% Not needed with XeLaTeX, which is needed to use the Erewhon font.
%\usepackage[utf8]{inputenc}

% Package that makes latex aware of the languages used in the text
\usepackage[brazil,english]{babel}

% Output font encoding for international characters
\usepackage[T1]{fontenc}

% Customize the header and footer
\usepackage{fancyhdr}

% Required for calculating the number of pages in the document
\usepackage{lastpage}

% Control sequences
\usepackage{etoolbox}

% Colors for links, text and headings
\usepackage[pdfborder={0 0 0}]{hyperref}

% Required for specification of custom fonts
\usepackage{fontspec}
\defaultfontfeatures{Ligatures=TeX}
\setmainfont[Path = ./fonts/,
Extension = .otf,
BoldFont = Erewhon-Bold,
ItalicFont = Erewhon-Italic,
BoldItalicFont = Erewhon-BoldItalic,
SmallCapsFeatures = {Letters = SmallCaps}
]{Erewhon-Regular}

 % Required for custom colors
%\usepackage{color}
%\definecolor{slateblue}{rgb}{0.17,0.22,0.34}

%\usepackage{sectsty} % Allows customization of titles
%\sectionfont{\color{slateblue}} % Color section titles

% Suppress section numbering
\setcounter{secnumdepth}{0}

% Define a custom page style
\fancypagestyle{plain}{\fancyhf{}\cfoot{\thepage\ of \pageref{LastPage}}}
% Use the custom page style through the document
\pagestyle{plain}
% Disable the default header rule
\renewcommand{\headrulewidth}{0pt}
 % Disable the default footer rule
\renewcommand{\footrulewidth}{0pt}

 % Stop paragraph indentation
\setlength\parindent{0pt}

% Non-indenting itemize
\newenvironment{itemize-noindent}
{\setlength{\leftmargini}{1em}\begin{itemize}}
{\end{itemize}}

% Text width for tabbing environments
\newlength{\tabwidth}
\setlength{\tabwidth}{3.5cm}
\newlength{\smallertextwidth}
\setlength{\smallertextwidth}{\textwidth}
\addtolength{\smallertextwidth}{-\tabwidth}

% Custom square bullet point definition
\newcommand{\sqbullet}{~\vrule height 1ex width .8ex depth -.2ex}

%----------------------------------------------------------------------------------------
%	MAIN HEADER COMMAND
%----------------------------------------------------------------------------------------

\renewcommand{\title}[1]{
{\huge{\textbf{#1}}}\\ % Header section name and color
\rule{\textwidth}{0.5mm}\\ % Rule under the header
}

%----------------------------------------------------------------------------------------
%	JOB COMMAND
%----------------------------------------------------------------------------------------

\newcommand{\employer}[4]{
\begin{tabbing}
\hspace{\tabwidth} \= \kill
\textbf{#1 - #2} \> \href{#4}{#3}
\end{tabbing}
}

\newcommand{\job}[2]{
\begin{tabbing}
\hspace{\tabwidth} \= \kill
\>\+ \textit{#1} \\
\begin{minipage}{\smallertextwidth}
\vspace{2mm}
#2
\end{minipage}
\vspace{2mm}
\end{tabbing}
}

%----------------------------------------------------------------------------------------
%	EDUCATION COMMAND
%----------------------------------------------------------------------------------------

\newcommand{\education}[6]{
\begin{tabbing}
\hspace{\tabwidth} \= \kill
\textbf{#1 - #2} \>\+ #3 - \href{#5}{#4} \\
\ifstrempty{#6}{}{
\begin{minipage}{\smallertextwidth}
\vspace{2mm}
#6
\vspace{2mm}
\end{minipage}
}
\end{tabbing}
}
%----------------------------------------------------------------------------------------
%	SKILL GROUP COMMAND
%----------------------------------------------------------------------------------------

\newcommand{\skillgroup}[2]{
\begin{tabbing}
\hspace{5mm} \= \kill
\sqbullet \>\+ \textbf{#1} \\
\begin{minipage}{\smallertextwidth}
\vspace{2mm}
#2
\end{minipage}
\end{tabbing}
}

%----------------------------------------------------------------------------------------
%	INTERESTS GROUP COMMAND
%-----------------------------------------------------------------------------------------

\newcommand{\interestsgroup}[1]{
\begin{tabbing}
\hspace{5mm} \= \kill
#1
\end{tabbing}
\vspace{-10mm}
}

% Define a custom command for individual interests
\newcommand{\interest}[1]{\sqbullet \> \textbf{#1}\\[3pt]}

%----------------------------------------------------------------------------------------
%	TABBED BLOCK COMMAND
%----------------------------------------------------------------------------------------

\newcommand{\tabbedblock}[1]{
\begin{tabbing}
\hspace{2cm} \= \hspace{4cm} \= \kill
#1
\end{tabbing}
}
 % Include the file specifying document layout

\hyphenation{li-de-ro}

%----------------------------------------------------------------------------------------

\begin{document}

%----------------------------------------------------------------------------------------
%	NAME AND CONTACT INFORMATION
%----------------------------------------------------------------------------------------

\title{\href{https://br.linkedin.com/in/laggarcia}{Leonardo Garcia}} % Print the main header

%------------------------------------------------

\parbox{0.5\textwidth}{ % First block
\begin{tabbing} % Enables tabbing
\hspace{3.5cm} \= \hspace{4cm} \= \kill % Spacing within the block
{\bf Data de nascimento} \> 28 de março de 1981 \\ % Date of birth 
{\bf Nacionalidade} \> Brasileira \\ % Nationality
{\bf Telefone residencial} \> +55 (19) 3241-5362 \\ % Home phone
{\bf Telefone celular} \> +55 (19) 99194-6670 \\ % Mobile phone
{\bf E-mail} \> \href{mailto:laggarcia@gmail.com}{laggarcia@gmail.com} \\ % Email address
\end{tabbing}}
\hfill % Horizontal space between the two blocks
\parbox{0.5\textwidth}{ % Second block
\begin{tabbing} % Enables tabbing
\hspace{3cm} \= \hspace{4cm} \= \kill % Spacing within the block
{\bf Endereço} \\
Rua Amador Bueno, 225 - torre 2 - ap. 44 - Vila Industrial \\ % Address line 1
Campinas - SP - Brasil \\ % Address line 2
CEP: 13035-030 \\ % Address line 3
\end{tabbing}}

%----------------------------------------------------------------------------------------
%	PERSONAL PROFILE
%----------------------------------------------------------------------------------------

\section{Perfil profissional}

% Leadership and development job application
Engenheiro de computação com mais de 15 anos de experiência em diversas áreas, procurando por uma posição desafiadora em uma empresa na qual possa contribuir não apenas com minhas habilidades técnicas, mas também com minhas habilidades de liderança. Gosto de montar e cultivar times, mantendo-os motivados mesmo sob situações adversas.

% Research job application
%A computer engineer with 15+ years of experience in a variety of areas, seeking a challenging and exciting position in a multi-national, multi-cultural, world wide research team working with High Performance Computing and Parallel Computing. I am good at building teams and working in tight collaboration with others.

%----------------------------------------------------------------------------------------
%	EDUCATION SECTION
%----------------------------------------------------------------------------------------

\section{Formação}

\education
{jan 2007}{nov 2013}
{Mestrado em Ciência da Computação}
{Universidade Estadual de Campinas (Unicamp)}{http://www.ic.unicamp.br/en}
{\href{http://repositorio.unicamp.br/jspui/handle/REPOSIP/275619}{Análise e estudo de desempenho e consumo de energia de memórias transacionais em software}}

%------------------------------------------------

\education
{mar 1999}{nov 2003}
{Engenharia de Computação}
{Universidade Estadual de Campinas (Unicamp)}{http://www.unicamp.br/unicamp/?language=en}


%----------------------------------------------------------------------------------------
%	EMPLOYMENT HISTORY SECTION
%----------------------------------------------------------------------------------------

\section{Experiência profissional}

\employer
{abr 2006}{presente}
{IBM Linux Technology Center (LTC)}{http://www.ibm.com}
\job
{Líder técnico de desenvolvimento em virtualização aberta e KVM em Power (jan 2013 - presente)}
{

\begin{itemize}
\item{Líder técnico de desenvolvimento de KVM em Power (ago 2014 - presente): lidero um time espalhado em 5 países e 8 fusos horários trabalhando com a comunidade de \href{http://www.qemu.org/}{QEMU}/\href{http://www.linux-kvm.org}{KVM}, \href{https://github.com/kimchi-project/kimchi}{Kimchi} e \href{https://github.com/kimchi-project/ginger}{Ginger} para Power, assim como desenvolvendo e mantendo estes projetos em \href{http://www-03.ibm.com/systems/power/software/linux/powerkvm/}{PowerKVM} e \href{https://open-power-host-os.github.io/}{OpenPOWER Host OS}.}
\item{Líder técnico de desenvolvimento em virtualização aberta (jan 2013 - ago 2014): liderei um time espalhado por EUA, Brasil e China trabalhando com a comunidade de \href{http://www.qemu.org/}{QEMU}/\href{http://www.linux-kvm.org}{KVM} e \href{https://github.com/kimchi-project/kimchi}{Kimchi} para x86 e código de virtualização aberta genérico para diversas plataformas.}
\end{itemize}

}
\job
{Líder técnico e desenvolvedor de distribuição Linux com virtualização aberta (jan 2010 - dez 2012)}
{Customização, automação de compilação e suporte para múltiplos clientes usando uma distribuição Linux mínima com suporte a virtualização.

\begin{itemize}
\item {Desenvolvi a configuração automática do hardware de máquinas virtuais.}
\item{Desenvolvi ferramentas para realizar o backup e a restauração de configurações do servidor.}
\item{Desenvolvi um cliente VNC que executa sobre um Linux framebuffer.}
\item{Tornei-me o líder técnico deste time em janeiro de 2011. Expandi o time para 3 países (Brasil, EUA e China).}
\item{Fui promovido a Engenheiro de Software Sênior enquanto neste projeto. Recebi um IBM Outstanding Technical Achievement Award (OTAA), uma das mais importantes distinções técnicas dentro da IBM, pelo meu trabalho neste projeto.}
\end{itemize}

}
\job
{Engenheiro de Software (jan 2009 - dez 2009)}
{

\begin{itemize}
\item{ServerGuide Scripting Toolkit, Linux Edition (mar 2009 - dez 2009): desenvolvi uma ferramenta que instalava e configurava automaticamente RHEL, SLES e VMware ESX em servidores Linux x86 e Power.}
\item{Shell Script Translator (jan 2009 - mar 2009): desenvolvi um analisador de sintaxe para ksh.}
\end{itemize}

}
\job
{Líder técnico e desenvolvedor do Cell IDE (abr 2006 - jan 2009)}
{Desenvolvimento de uma IDE (ambiente integrado de desenvolvimento) C/C++ baseada em \href{https://eclipse.org/}{Eclipse} para os processadores da IBM Cell/B.E. e wire-speed Power.

\begin{itemize}
\item{Tornei-me o líder técnico deste time em outubro de 2007.}
\item{Integrei as ferramentas de compilação e análise de desempenho; desenvolvi um gerador automático de makefiles.}
\item{Desenvolvi um framework para execução remota.}
\item{Responsável pela compilação do Cell IDE.}
\item{Interação com as comunidades do Eclipse e de seus plug-ins; interação com as comunidades de usuários.}
\item{Contribui o Cell IDE para o \href{http://www.eclipse.org/ptp/}{PTP}.}
\end{itemize}

}
\job
{Outras atividades}
{Na IBM, também me envolvi com:

\begin{itemize}
\item{Gerente de projeto e consultor de P\&D (mar 2009 - presente): projetos desenvolvidos em parceria com institutos de P\&D externos através de programas de isenção de impostos do governo brasileiro. Trabalhei como gerente de projeto nos seguintes projetos:

\begin{itemize}
\item{Microsoft Active Directory - Samba 4 migration (mar 2009 - abr 2010).}
\item{Virtualização aberta em Power Systems (mar 2010 - mar 2011): VirtIO em PowerVM.}
\end{itemize}

}
\item{Recrutador técnico (jan 2008 - presente).}
\item{Consultor técnico (jan 2008 - presente): análise e ajuda com problemas no ecossistema Linux relacionados a pré-vendas, pós-vendas ou situações de suporte vindas de outras áreas da companhia.}
\item{Programa de mentorização (jan 2007 - presente).}
\item{Embaixador universitário (set 2006 - jan 2010).}
\end{itemize}

}

%------------------------------------------------

\employer
{mar 2003}{abr 2006}
{Compera (atual Movile)}{https://www.movile.com/en/}
\job
{Analista de suporte a software, operações e desenvolvimento}
{Trabalhei em múltiplas frentes do Dispara, uma plataforma de gerência de forças de campo que integra, através de dispositivos móveis, o back-office das empresas com suas forças de campo. Comecei trabalhando com suporte e monitoramento da plataforma dentro de um cliente e cresci minhas responsabilidades até o ponto em que era responsável por todo o produto.

\begin{itemize}
\item{Gerente de suporte e monitoramento: suporte 24x7 e monitoramento de todos os clientes do Dispara: mais de 4.500 usuários por todo o Brasil trabalhando em mais de 40 empresas.}
\item{Analista de operações: responsável pela operação, manutenção, disponibilidade e desempenho da plataforma Dispara, incluindo o tuning do banco de dados SQL Server.}
\item{Analista de projetos: responsável pela análise dos requisitos dos clientes.}
\item{Analista de desenvolvimento: líder do time de especificação técnica e desenvolvimento de integrações; responsável pelas atividades de projeto, desenvolvimento e implantação.}
\item{Gerente de produto: responsável pela estrutura, modelo e desenvolvimento das novas versões da plataforma Dispara; integrei o Dispara com múltiplas telecoms (via SMS, UP.Notify e WAP Push); ajudei no desenho de um gerente de fluxo de trabalho}
\end{itemize}

}

%------------------------------------------------

\employer
{fev 2003}{mar 2003}
{Dextra}{http://dextra.com.br/en/}
\job
{Estagiário de desenvolvimento de software}
{Desenvolvi web sites dinâmicos baseados em Java (JSP e Servlets) e Perl.}

%------------------------------------------------

\employer
{dez 2001}{mar 2003}
{CenPRA}{http://www.cenpra.gov.br/}
\job
{Estagiário de pesquisa}
{Trabalhei com micro-controladores, módulos de Linux real-time e uma distribuição Linux para o Autonomous Unmanned Remote mOnitoring Robotic Airship (AURORA).

\begin{itemize}
\item{Desenvolvi o hardware e o software em micro-controladores PIC para obter medidas feitas por sensores e enviá-las através de uma Control Area Network (CAN).}
\item{Integrei o sistema de visão desenvolvido para o dirigível e seu sistema operacional (Real-time Linux).}
\item{Construi uma distribuição Linux do zero baseada em Slackware que atendia os requisitos de baixo consumo de memória e suporte a tempo-real.}
\item{Portei os módulos de Linux real-time do Linux 2.0 para o Linux 2.4.}
\item{Ajustei o simulador do dirigível para a nova plataforma de desenvolvimento e integrei o sistema de visão computacional nele.}
\end{itemize}

}

%----------------------------------------------------------------------------------------
%	PUBLICATIONS SECTION
%----------------------------------------------------------------------------------------

\section{Publicações}

\begin{itemize}
 \item \textit{Arquitetura de Computadores: educação, ensino e aprendizado - Capítulo 5: Proposta de um Elenco de Disciplinas de Pós-Graduação para Formação de Arquitetura de Computadores.} \\
  S. T. Kofuji ; J. M. Kofuji ; M. L. Netto ; W. J. Chau ; F. Valiante Filho ; \textbf{L. A. G. Garcia} ; E. D. Moreno ; E. Horta ; E. Lima ; D. Vahia.  \\
  ISBN: 978-85-7669-263-8, 2012, Sociedade Brasileira de Computação (SBC).
 \item \textit{Energy-Performance Tradeoffs in Software Transactional Memory.} \\
  A. Baldassin ; J. P. L. de Carvalho ; \textbf{L. A. G. Garcia}, Rodolfo Azevedo. \\
  Proceedings of the 24th International Symposium on Computer Architecture and High Performance Computing (SBAC-PAD'2012).  
 \item \textit{A Software Transactional Memory System for an Asymmetric Processor Architecture.} \\
  F. Goldstein ; A. Baldassin ; P. Centoducatte ; R. Azevedo ; \textbf{L. A. G. Garcia}. \\
  Proceedings of the 20th International Symposium on Computer Architecture and High Performance Computing (SBAC-PAD'08), v. 1, p. 175-182, 2008. \\
  Este artigo recebeu o prêmio Júlio Salek Aude de melhor artigo do SBAC-PAD'08.
\end{itemize}

%----------------------------------------------------------------------------------------
%	PRESENTATIONS SECTION
%----------------------------------------------------------------------------------------

\section{Apresentações}

\begin{itemize}
 \item \textit{Análise e estudo de desempenho e consumo de energia de memórias transacionais em software} \\
  \textbf{L. A. G. Garcia}. \\
  Melhor dissertação de mestrado do Concurso de Teses e Dissertação do XV Workshop em Sistemas Computacionais de Alto Desempenho (WSCAD 2014).
 \item \textit{Análise de desempenho e consumo de energia em Memórias Transacionais.} \\
  \textbf{L. A. G. Garcia} ; A. Baldassin ; R. Azevedo. \\
  Poster apresentado na III Escola Regional de Alto Desempenho de São Paulo (ERAD-SP 2012).
 \item \textit{FaceID: Benchmarks for Face Detection Algorithms Using Cell Processor.} \\
  R. K. Hiramatsu ; J. M. Kofuji ; S. T. Kojufi ; L. Mattes ; \textbf{L. A. G. Garcia}. \\
  Poster na The 2nd Annual Carleton Cell BE Programming Workshop, 2009.
 \item \textit{Programando multicore com IBM full-system simulator - Cell Broadband Engine.} \\
  J. M. Kofuji ; S. T. Kojufi ; R. K. Hiramatsu ; \textbf{L. A. G. Garcia}. \\
  Apresentação no VIII Workshop em Sistemas Computacionais de Alto Desempenho (WSCAD 2007).
 \item \textit{Workshop de Programação Cell}. \\
  \textbf{L. A. G. Garcia}. \\
  Workshop no Congresso Regional de Computação e Informática do Oeste Paulista, 2007.
\end{itemize}

%----------------------------------------------------------------------------------------
%	COMPUTING SKILLS SECTION
%----------------------------------------------------------------------------------------

\section{Habilidades em engenharia de software}

\skillgroup{Linguagens de programação}
{
\textit{C/C++, Java, C\#, Pascal, Assembly (x86, MIPS, Motorola 68k e Microchip PIC).} \\
\textit{Python, Perl, Bash, JavaScript, VBScript.} \\
}

%------------------------------------------------

\skillgroup{Outras}
{
\textit{Experiência com vários sistemas operacionais}: Linux, Unix (Solaris, Digital), Windows. \\
\textit{Arquitetura e administração de base de dados}: SQL Server, Firebird. \\
\textit{Programação para redes}: protocolos TCP/IP (HTTP, FTP, TELNET, POP, SMTP), programação cliente-servidor, sockets, RPC, RMI, CORBA e Web Services. \\
\textit{Conhecimento sobre metodologias de desenvolvimento tradicionais e Ágeis.} \\
\textit{Controle de versão de conteúdo}: git, SVN, CVS. \\
}

%----------------------------------------------------------------------------------------
%	LANGUAGES
%----------------------------------------------------------------------------------------

\section{Línguas}

\textit{Português nativo} \\
\textit{Inglês proficiente} \\
\textit{Espanhol e francês: entendimento básico} \\

\end{document}
